\clearpage
\section[Human Eval Results]{Human Eval Results \hyperref[chsevenrq2]{\chsevenrqtwo}}\label{sec:ch-7-human-eval-results}

We collected 140 completed query sessions from 20 annotators, yielding 420 pairwise comparisons across the three-match tournament described in~\Cref{sec:ch-7-human-eval-setup}. On average, annotators spent roughly four minutes per page (median 258s). Sixteen annotators completed all 8 assigned queries.

\paragraph{DeepSeek-V3.1 is preferred over Llama-3.3-70B.}
Aggregating across both prompting conditions and all tournament rounds, DeepSeek-V3.1 wins 72.8\% of its matchups (91/125) in the final round, against 27.2\% for Llama-3.3-70B-Instruct. The gap is even more pronounced in direct cross-model first-round matchups, where DeepSeek wins 73.1\% of head-to-head comparisons (68/93). This preference is consistent across the four-variant breakdown shown in \Cref{tab:variant-winrates}. Both DeepSeek variants outperform both Llama variants in both rounds.

\paragraph{Personalization is preferred when generated by a stronger model.}
The effect of conditioning on the persona profile depends on where in the tournament the comparison occurs (\Cref{tab:personalization}). In the final round, annotators prefer the personalized summary 63.2\% of the time (79/125) over the generic summary, with 15 ties (``neither''). The first-round picture is less clean: in the \emph{A vs.\ B} bracket, personalized summaries win only 46.5\% of the time (60/129) while in the \emph{C vs.\ D} bracket, personalized summaries win 62.2\% (79/127), matching the final-round rate. We hypothesize this effect is a result of the fact that the first-round result aggregates across both models, and our cross-model comparisons are dominated by a strong preference for DeepSeek-V3.1 over Llama-3.3-70B (\Cref{tab:variant-winrates}). If Llama is simply less effective at personalization, then including its personalized outputs in the first-round average would dilute the personalization signal there, while the final round, which is disproportionately populated by the stronger model's outputs, would surface it more cleanly. 

\begin{table*}[h!]
\centering
\begin{subtable}[t]{0.75\textwidth}
\centering
\footnotesize
\begin{tabular}{lcccc}
& \multicolumn{2}{c}{\textbf{Round 1}} & \multicolumn{2}{c}{\textbf{Final}} \\
\cmidrule(lr){2-3}\cmidrule(lr){4-5}
\textbf{Variant} & W/N & \% & W/N & \% \\
\midrule
DeepSeek-V3.1 (personalized) & 88/140 & 62.9 & 58/88 & \textbf{65.9} \\
DeepSeek-V3.1 (generic)      & 75/140 & 53.6 & 33/75 & 44.0 \\
Llama-3.3-70B (personalized) & 51/140 & 36.4 & 21/51 & 41.2 \\
Llama-3.3-70B (generic)      & 42/140 & 30.0 & 13/42 & 31.0 \\
\end{tabular}
\caption[Win rates for each of the four summary variants.]{Win rates for each of the four summary variants in round 1 (out of 140 round-1 matchups per variant) and in the final round (out of the variant's round-1 wins). DeepSeek-V3.1 with persona conditioning is the most preferred variant across both rounds.}
\label{tab:variant-winrates}
\end{subtable} 
\vspace{1em}
\begin{subtable}[t]{0.75\textwidth}
\centering
\footnotesize
\begin{tabular}{lccc}
 & & & \\
\textbf{Comparison} & \textbf{Pers.} & \textbf{Gen.} & \textbf{Tie} \\
\midrule
Round 1 (A vs.\ B) & 60 (46.5\%) & 69 (53.5\%) & 11 \\
Round 1 (C vs.\ D) & 79 (62.2\%) & 48 (37.8\%) & 13 \\
Final round        & \textbf{79 (63.2\%)} & 46 (36.8\%) & 15 \\
\hdashfull
Same-model R1 & 51 (57.3\%) & 38 (42.7\%) & --- \\
\end{tabular}
\caption[Annotator preferences between persona-conditioned and generic summaries.]{Annotator preferences between persona-conditioned and generic summaries, by tournament position. Personalization is preferred in the final round and in the second first-round bracket (C vs.\ D), and---when model identity is held fixed---in same-model first-round matchups.}
\label{tab:personalization}
\end{subtable}
\caption[Information satisfaction human evaluation results.]{Human evaluation results. (\subref{tab:variant-winrates}) Per-variant win rates across the $2 \times 2$ design. (\subref{tab:personalization}) Aggregated personalized-vs-generic preferences by tournament position.}
\label{tab:ch-7-human-eval}
\end{table*}



\paragraph{Automatic and LLM-based metrics correlate poorly with human judgment.}
To assess whether existing summarization metrics recover the human preferences just described, we compute the pairwise agreement rate between each metric and the annotator's choice on every comparison in our evaluation. A metric agrees with the annotator on a given pair if it assigns a higher score to the summary that the annotator selected; chance agreement is 50\%. We count the metric's decision as ``Neither'' if the scores are within a per-metric threshold of similarity. 

\Cref{tab:metric-agreement} reports agreement rates grouped into three families: reference-based and reference-free automatic metrics (e.g., ROUGE, BERTScore, BLANC), LLM-as-a-judge scores using Prometheus-7B as the evaluator, and LLM-as-a-judge scores using Llama-3.3-70B-Instruct as the evaluator. Across all three families, agreement rates cluster tightly around chance: automatic metrics range from 45\% to 60\% (mean 51.6\%), Prometheus-based judgments range from 46\% to 55\% (mean 50.8\%), and Llama-based judgments range from 47\% to 53\% (mean 47.7\%, excluding one prompt for which Llama failed to follow instructions). Notably, persona-aware variants we designed specifically to capture informational satisfaction---\textsc{persona\_precision} and \textsc{persona\_recall}---also fail to exceed chance (47\%--51\%).

The strongest performers are surface-overlap metrics anchored to a reference summary (ROUGE-2 F1 at 60\%, ROUGE-1 at 57\%, chrF at 58\%), but even these are far from the level required to use them as proxies for personalized informational satisfaction. The LLM-judge protocols, including those explicitly conditioned on the annotator's persona profile, perform no better than reference-free automatic metrics. Together, these results support the central claim of this chapter: existing summarization metrics, including state-of-the-art LLM-based judges, are insufficient measures of how well a summary serves an individual reader's informational needs.


\begin{table}[h!]
\centering
\footnotesize
\begin{subtable}[t]{0.44\textwidth}
\centering
\begin{tabular}{lc}
\toprule
\textbf{Metric} & \textbf{Agreement} \\
\midrule
ROUGE-1 F1        & 57\% \\
ROUGE-2 F1        & \textbf{60\%} \\
ROUGE-L F1        & 52\% \\
chrF              & 58\% \\
SummaQA F1        & 56\% \\
BERTScore F1      & 52\% \\
BLEU              & 53\% \\
METEOR            & 53\% \\
BLANC             & 50\% \\
SUPERT            & 50\% \\
CIDEr             & 50\% \\
\bottomrule
\end{tabular}
\caption{Traditional metrics}
\end{subtable}
\hfill
\begin{subtable}[t]{0.55\textwidth}
\centering
\begin{tabular}{lc}

\toprule
\textbf{Metric} & \textbf{Agreement} \\
\midrule
Overall                  & 49\% \\
Relevance                & 51\% \\
Coherence                & \textbf{53\%} \\
Consistency              & 50\% \\
Fluency                  & 46\% \\
Informativeness          & \textbf{53\%} \\
FactScore                & 51\% \\
Persona-precision        & 47\% \\
Persona-recall           & 51\% \\
Persona-relative overall & 50\% \\
\bottomrule
\end{tabular}
\caption{LLM-as-a-judge (Prometheus-7B)}
\end{subtable}
\vspace{1em} \\
\begin{subtable}[t]{0.55\textwidth}
\centering
\begin{tabular}{lc}
\toprule
\textbf{Metric} & \textbf{Agreement} \\
\midrule
Overall            & 51\% \\
Relevance          & 50\% \\
Coherence          & 51\% \\
Consistency        & 50\% \\
Informativeness    & 49\% \\
FactScore          & \textbf{52\%} \\
Persona-precision  & 49\% \\
Persona-recall     & 49\% \\
\bottomrule
\end{tabular}
\caption{LLM-as-a-judge (Llama-3.3-70B-Instruct)}
\end{subtable}
\caption[Pairwise agreement rate between metrics and human pairwise preferences.]{Pairwise agreement rate between metrics and human pairwise preferences. Chance agreement is 50\%. No metric family substantively exceeds chance, including persona-aware variants designed to capture informational satisfaction.}
\label{tab:metric-agreement}
\end{table}